\subsection{Character consequences and dependency repairs}


\begin{statement}{B.1}{Proved}{The precise finite-frame character classification}
Among continuous characters of $G=\GZ$ with values in $\Ftwo$
which vanish on $G_{\Q}$ and whose restriction to $\HH$ factors
through the two corrected finite frames $C_6\times C_4$, the
only characters are zero and $\Delta$.
\end{statement}
\begin{proof}
The full frame image and the Galois image are respectively
$C_6\times C_4$ and its index-two diagonal-parity subgroup,
as proved in B.14 below. A homomorphism from their quotient
to $\Ftwo$ is either zero or its unique nonzero character.
Thus the restriction is zero or $\delta$. The decomposition
$G=\HH G_{\Q}$ then uniquely determines a character
vanishing on $G_{\Q}$. This proves the stated restricted
classification, and makes no claim that every relative
character factors through these frames.
\end{proof}

\begin{statement}{B.2}{Proved}{The exact meaning of the constant sign function}
On $\HH$, let $\ell_h$ be the function on the two-element set
$\{j=0,j=1728\}$ with values
$((-1)^{a_3(h)},(-1)^{a_2(h)})$. Every such function occurs,
and it is constant exactly on $\ker\delta$.
On all of $G$, the gauge-invariant version is its class modulo
the common sign; its zero locus is exactly $\ker\Delta$.
This is the largest closed subgroup on which that obstruction
vanishes. No analogous statement about all five-point special
points follows without constructing their coordinates.
\end{statement}
\begin{proof}
R2 gives the pair $(1,0)$, while A.3 supplies a Galois pair
$(1,1)$. The image of the homomorphism $(a_3,a_2)$ therefore
contains all four elements of $\Ftwo^2$. Equality of the two
signs is the equation $a_3+a_2=0$. A subgroup on which the
equation holds is contained in its kernel, which is closed
and itself has the property. A.2--A.3 prove the asserted
all-multiplier formulation and explain why its two separate
components cannot be extended as abelian characters.
\end{proof}

\begin{statement}{B.3}{Proved}{A tower-action module carrying the sign line}
There is a finite-rank free $\Z_\ell$-module, for every odd
prime $\ell$, carrying a continuous $G$ action through its
action on a finite characteristic quotient of $\B$, and
containing a line with character $(-1)^\Delta$.
This construction is a permutation module of finite
automorphism residues; it is not a proof that a specified
level-12 homology module has such a line.
\end{statement}
\begin{proof}
Choose the finite characteristic level of A.10, let $Q$
be the finite image of $G$ in its outer automorphism group,
and let $\bar\Delta:Q\to\Ftwo$ be the induced character.
On $M=\Z_\ell[Q]$ use left translation of the basis $e_q$.
Then
\[
 v=\sum_{q\in Q}(-1)^{\bar\Delta(q)}e_q
\]
is nonzero, primitive, and satisfies
$g v=(-1)^{\Delta(g)}v$. The equality follows by replacing
$q$ with $g^{-1}q$ in the sum. The line is a direct
$\Z_\ell$-summand because one of its coordinates is a unit.
All maps factor through the specified finite tower action.
\end{proof}

\begin{statement}{B.4}{Refuted}{Not every module is distinguished from its sign twist}
Every $G$-module and its $\Delta$ twist restrict to canonically
identical $G_{\Q}$ actions. They need not be nonisomorphic
as $G$-modules, even when they contain the sign line.
\end{statement}
\begin{proof}
Let $\chi=(-1)^\Delta$ and take
$M=\Z_\ell\oplus\Z_\ell(\chi)$. The twisted module is
$\Z_\ell(\chi)\oplus\Z_\ell$, and interchanging the
summands is a $G$-equivariant isomorphism. On an element
with $\Delta=1$ the matrices are
\[
 D=\begin{pmatrix}1&0\\0&-1\end{pmatrix},\quad
 -D,\qquad
 P=\begin{pmatrix}0&1\\1&0\end{pmatrix},\quad PD=(-D)P.
\]
The executable algebra check verifies this identity.
On $G_{\Q}$ the twist is literally one, so restriction
does not change any action. In contrast the rank-one
trivial representation and its twist are nonisomorphic
over $\Z_\ell$ for odd $\ell$, since a putative
isomorphism would intertwine $1$ and $-1$ on the R2
element. This is the precise true replacement for the
universal nonisomorphism assertion.
\end{proof}

\begin{statement}{B.5}{Proved}{The powerset semiring statement, with its operations specified}
Let $\mathscr C(G)$ be the set of closed subsets of the
profinite group $G$, with addition given by union and
multiplication by setwise group product. Then
$S\mapsto\Delta(S)$ is a surjective homomorphism of
idempotent semirings to the powerset semiring of $\Ftwo$.
The image from closed subsets of $G_{\Q}$ is
$\{\varnothing,\{0\}\}$. The same holds for $\GT$
if and only if $\Delta|_{\GT}=0$.
\end{statement}
\begin{proof}
Finite unions of closed sets are closed, and products
are compact images of products of compact spaces, hence
closed. Group multiplication distributes over unions.
The zero is $\varnothing$ and the unit is $\{1\}$.
Direct image preserves unions, and the homomorphism law
gives $\Delta(ST)=\Delta(S)+\Delta(T)$ as sumsets.
The empty set, singleton identity, singleton R2 element,
and their two-element union realize all four subsets of
$\Ftwo$. Galois vanishing proves its asserted image.
The last equivalence follows by using singleton subsets.
If ``semiring'' were instead meant to use intersection
as multiplication, the assertion would be false: the
disjoint sets $\{1\}$ and $\{c\}$ both have image
$\{0\}$. Thus the operations matter.
\end{proof}

\begin{statement}{B.6}{Proved}{Continuous ring maps are precisely continuous unit characters}
For odd $\ell$, continuous unital ring homomorphisms
\[
 \Z_\ell[[G]]\longrightarrow\Z_\ell
\]
which agree with augmentation on $\Z_\ell[[G_{\Q}]]$
are in bijection with continuous characters
$\chi:G\to\Z_\ell^\times$ trivial on $G_{\Q}$.
If $K_G$ is the relative abelianization in P.12,
this set is
\[
 \Hom_{\rm cont}(K_G,\mu_{\ell-1})
 \times\Hom_{\rm cont}(K_G,\Z_\ell).
\]
In particular augmentation and $g\mapsto(-1)^{\Delta(g)}$
are distinct such maps. This is an exact classification
in terms of $K_G$, whose independent computation is
still required to list all characters explicitly.
\end{statement}
\begin{proof}
A unital continuous map fixes the ordinary integers
and therefore fixes $\Z_\ell$ by density. Group
elements are units and give a continuous character.
Conversely, reduction of a continuous unit character
modulo $\ell^n$ has finite image and open kernel.
Extending linearly on the corresponding finite group
ring and passing to the inverse limit gives the ring
map; density of the group algebra proves uniqueness.
Triviality on $G_{\Q}$ is precisely the augmentation
condition. Factorization through the abelianization
and its split cyclotomic factor gives $K_G$.
Finally
$\Z_\ell^\times=\mu_{\ell-1}\times(1+\ell\Z_\ell)$
and the convergent $\ell$-adic logarithm identifies
the second factor with the additive group $\ell\Z_\ell$,
which is isomorphic to $\Z_\ell$. The R2 value makes
the two displayed maps distinct.
\end{proof}

\begin{statement}{B.7}{Proved}{The relative residual tower detects every element}
For cofinal characteristic open subgroups $U$ of $\B$,
\[
 \B\simeq\varprojlim_U\B/R_{\rm nil}(U).
\]
An automorphism compatible with every quotient and
every transition map is therefore an actual automorphism
of $\B$. Reaching each quotient separately does not
produce a compatible system of automorphisms satisfying
an additional equation such as the pentagon.
\end{statement}
\begin{proof}
If $V\subset U$, restriction of a finite nilpotent
quotient of $U$ shows
$R_{\rm nil}(V)\subset R_{\rm nil}(U)$.
Also $R_{\rm nil}(U)\subset U$, so the intersection
of the residuals is trivial. A finite family of coset
conditions on the right can be combined in a finer
quotient. Compactness gives surjectivity and the
intersection calculation gives injectivity. Compatible
automorphisms and their compatible inverses act on
this inverse limit. This states the required transition
compatibility rather than assuming that arbitrary
members of the individual automorphism groups form
an inverse system. Requiring a continuous word
identity at every level is another closed condition;
its solvability is not implied by independent lifts.
\end{proof}

\begin{statement}{B.8}{Proved}{The primitive, the cocycle, and the numerical coordinate are distinct}
In the normalized multiplier-one modular action,
\[
 \Phi(r)=ArA^{-1}=f^{-1}r
 \quad\Longrightarrow\quad
 f^{-1}=A\rho(A)^{-1},\qquad \rho=\operatorname{Ad}(r)|_F.
\]
Thus the manuscript's $g$ in
$u\rho=\operatorname{int}(g)\rho u$ is a cocycle,
and $A$ is its primitive. It is not the same object
as the Gaussian primitive $\beta$, nor as the scalar
$\rho_\Phi$ defined by
$(\beta^{-1})^{\rm ab}=(\rho_\Phi,\rho_\Phi)$.
\end{statement}
\begin{proof}
Multiply the displayed conjugation equation on the
right by $r^{-1}$. This gives
$f^{-1}=ArA^{-1}r^{-1}=A\rho(A)^{-1}$.
The three objects have different types: $g,A,\beta$
are group elements, whereas $\rho_\Phi$ is a profinite
integer. The cocycle and a primitive cannot be
interchanged in a finite evaluation. Cote 145,
scans 7, 11--12, 22 and 28, is available here only
through the supplied first-pass transcription;
the exact dictionary for the manuscript's separate
letter $\beta$ requires its written relation and
cannot be inferred merely from a modern reuse of
the same letter. B.24--B.25 below now compute it from the supplied scan-12 relation,
and prove that it is not $\rho_\Phi$.
\end{proof}

\begin{statement}{B.9}{Proved}{The normalized pair determines the outer gentle element}
In the fixed free generators and Belyi normalization,
the pair $(\lambda,f)$ with $f\in F'$ determines
an outer gentle element injectively. This statement
does not say that its residue modulo 12 determines
the two primitive frames.
\end{statement}
\begin{proof}
The pair specifies $x\mapsto x^\lambda$ and
$y\mapsto f^{-1}y^\lambda f$, hence the automorphism.
Conversely two such normalized representatives of
the same outer action differ by conjugation with
an element centralizing $x^\lambda$, hence by a
power $x^a$. Comparing the $y$ images shows that
the two conjugators differ by a power of $y$ and
this power of $x$. Their zero abelianizations
force both powers to vanish; centralizers of
nontrivial generator powers in a free profinite
group are the corresponding procyclic subgroups.
Thus the normalized $f$ is unique. A.9 gives a
specific reason why a small finite action can lose
the primitive-frame information.
\end{proof}

\begin{statement}{B.10}{Proved}{The corrected dependency order used in this return}
Every reference here to the CM field statement is
written ``LM7 D.25'', rather than the ambiguous
``D.25''. The R4 input is the corrected E.65 proof
including the pure-frame existence argument and
the explicit two-model Galois sign comparison.
The lifting and weight proofs in Part P and B.11--B.13
are ordered without the circular dependencies
E.2/E.58, D.1/D.51 and D.2/D.52.
\end{statement}
\begin{proof}
This is a statement about the dependencies of this
text, which can be checked from its displayed
proofs. The primitive existence uses the torsion
conjugacy and centralizer theorem, specifically
Herfort--Ribes (1985), Theorem B$'$, and the
congruence convention uses LM7 D.15. The explicit
field calculation appears independently in A.13;
E.62 is not an input. The positive residue in R2
is obtained from the relative lifting theorem
before it is used to prove that the sign is nonzero.
The weight computations are integrated before
their consequences are stated. The Hopficity
argument is the finite-open-subgroup permutation
argument in P.7, not an implication from finite
surjectivity without injectivity. No use of Turaev
2009 is made. These repairs concern the present
paper and its register; they do not silently
rewrite the supplied archive.
\end{proof}

\subsection{Weights, the complete frame image and the arithmetic replacement}


\begin{statement}{B.11}{Proved}{the full special weight-five image}
\end{statement}
For every prime $\ell$, the rational weight-five metabelian leading
image of the multiplier-one full peripheral outer automorphism
group is
\[
 (X+Y)\langle X^2,XY,Y^2\rangle_{\mathbf Q_\ell}[x,y].
\]
It has dimension three. Thus LM7\_PUSH D.51 is correct, and
its rational dimension statement also extends to two and three.

\emph{Proof.}
In the special equation $[h,y]=[k,x+y]$, take the actual quotient
of the free Lie algebra by its second derived ideal. If the
coefficients are $p,q$, then $Yp=(X+Y)q$. Since $Y,X+Y$ are
relatively prime, $X+Y$ divides the homogeneous cubic $p$.
This proves the upper bound three.

For the lower bound the three pairs in LM7\_PUSH D.51.3 were
expanded anew in \texttt{check\_lie\_witnesses.py}. All three
special equations vanish coefficientwise. Right-$y$ Fox
extraction, followed by abelianization, gives $XY$ times their
metabelian coefficient. The output is respectively
$(X+Y)X^2,(X+Y)XY,(X+Y)Y^2$. Their only possible denominator is
two. For odd $\ell$, P.7 integrates the displayed vectors in
the actual free pro-$\ell$ group. At $\ell=2$, integrate twice
each displayed pair; their rational images still span the same
three-dimensional $\mathbf Q_2$-space. No coprime $S_3$ step is
needed for this unrestricted peripheral statement.
Place each automorphism in the $\ell$ factor of
$F/R_{\rm nil}(F)=\prod_pF^{(p)}$, and place the identity in all
other factors, with the same tuple of peripheral conjugators.
The free-group assertion of R1 lifts this exact automorphism.
The characteristic pro-$\ell$ projection and the lower-central
and metabelian quotient maps retain the three leading terms.
Their linear independence proves the lower bound.

The representation-theoretic upper bound in D.21 also follows:
division by $X+Y$ identifies the target with quadratics in the
standard two-dimensional $S_3$ representation. In characteristic
zero this is the invariant line $X^2+XY+Y^2$ plus an irreducible
two-dimensional summand. An $S_3$-stable image containing the
invariant line therefore has dimension one or three; the actual
source construction decides it is three. \hfill$\square$

\begin{statement}{B.12}{Proved}{the full gentle weight-nine image}
\end{statement}
For every $\ell\geq5$, the rational weight-nine metabelian leading
image of $\widehat{\mathrm{GT}}_0$ is
\[
 \mathbf Q_\ell G_9[x,y]\oplus\mathbf Q_\ell P_9[x,y],
 \qquad G_9=\frac{(X+Y)^9-X^9-Y^9}{XY},\quad
 P_9=X^2Y^2(X+Y)^3.
\]
There is an actual normalized full pair whose first pro-$\ell$
coefficient has metabelian image $P_9[x,y]$.

\emph{Proof.}
The 42 displayed coefficients of $2h$ in LM7\_PUSH D.52.3
were entered afresh and expanded over $\mathbf Q$. The
antisymmetry, cabling, special and cyclic equations all vanish
coefficientwise; the associative support has 412 terms. Its Fox
image is exactly $XY P_9$, with coefficient vector $(1,3,3,1)$
in the four indicated monomials. Hence this is an integral
invariant special source at every $\ell\geq5$. P.7--P.9 and
R1 give the actual full pair with this leading term and its
correct positive pair sign. The derivation coordinate is its
negative, because $y\mapsto f^{-1}yf$.

For the upper bound let $p$ be the homogeneous degree-seven
metabelian coefficient of an invariant special leading term.
Its equations include
\[
 p(X,Y)=p(Y,X),\qquad
 Xp(X,Y)+(X+Y)p(-X-Y,Y)=0,
\]
and its cyclic sum is zero. Comparing coefficients gives
\[
\begin{split}
p={}&A(X^7+Y^7)+4A(X^6Y+XY^6)\\
   &+B(X^5Y^2+X^2Y^5)+(3B-14A)(X^4Y^3+X^3Y^4)\\
 ={}&\frac A9G_9+\left(B-\frac{28A}{3}\right)P_9.
\end{split}
\]
Thus the upper bound is two. We also independently formed all
56 Lyndon basis vectors in degree nine, formed the 2002 by 56
integer matrix of the three identities, and computed its exact
rational kernel. It has dimension four and metabelian rank two.
An explicit invariant free-Lie preimage of $G_9$ with denominator
four is retained in \texttt{invariant\_dimensions.json}. P.7--P.9
and R1 integrate that vector as well. This gives both directions
without using a circular dependence on D.1 or D.2, or an unproved
Galois nonvanishing statement. The exact matrix and all kernel
vectors are reproducible from \texttt{check\_invariant\_dimensions.py}.
\hfill$\square$

\begin{statement}{B.13}{Known}{every weight-nine lift with this metabelian coefficient fails}
\citep[Theorem 0.1]{LiftEnriquez2006}
\end{statement}
Enriquez's Theorem 0.1 states that a homogeneous degree-$m$
element of $\mathfrak{grt}_1(k)$ has metabelian image proportional
to $G_m[x,y]$ for odd $m\geq3$, and zero in other degrees.
Here $k$ is any characteristic-zero field. Consequently every
full gentle pair with no pro-$\ell$ terms below degree nine and
with degree-nine metabelian image $P_9[x,y]$ fails the pentagon,
regardless of its other degree-nine or higher coefficients.

\emph{Proof of the application.}
If a full pair in this stratum satisfied the group pentagon, its
first nonzero term would satisfy the infinitesimal pentagon, by
the continuous comparison and first-degree argument of P.11.
Its two group cycles likewise give the antisymmetry and cyclic
Lie identities. It therefore belongs to
$\mathfrak{grt}_1(\mathbf Q_\ell)_9$ and Enriquez applies.
The coefficient of $X^7$ in $G_9$ is 9, whereas it is zero in
$P_9$, and $P_9$ is nonzero. Thus they are not proportional.
The contradiction excludes every lift in this entire stratum.
This also supplies the requested pentagon upper bound of
LM7\_PUSH D.20; it does not by itself prove the separate
Galois lower-bound assertion in that row. \hfill$\square$

\begin{statement}{B.14}{Proved}{the finite image}
\end{statement}
In the signed corrected coordinates of R4,
\[
 (T,W):\mathcal H\longrightarrow C_6\times C_4
\]
is onto. The image of $G_{\mathbf Q^{\rm ab}}$ is precisely
\[
 J=\{(T,W):\epsilon_3(T)=\epsilon_2(W)\},
\]
which is cyclic of order 12.

\emph{Proof.}
Write $C_6=\langle -I\rangle\times\langle R_0\rangle$ and
$C_4=\langle S\rangle$, with the matrices and signed lift
which are identity modulo four as in R4. R4 gives the
homomorphisms and the containment of the Galois image in $J$.
Writing the two cyclic exponents as $i\bmod6,j\bmod4$ identifies
$J$ with $i\equiv j\pmod2$; it is generated by $(1,1)$ and has
order 12.

The three-division polynomial of $E_3:y^2=x^3+1$ is
$3x(x^3+4)$. The nonzero roots have $y^2=-3$, so
\[
 \mathbf Q(E_3[3])=\mathbf Q(\zeta_3,\sqrt[3]4).
\]
Its Galois group over $\mathbf Q$ is $S_3$, since $X^3-4$ is
irreducible and has nonsquare discriminant. Its intersection
with $\mathbf Q^{\rm ab}$ is the quadratic fixed field of
$A_3$, so the action over $\mathbf Q^{\rm ab}$ has order three.
The normalization here is essential. In the transported Legendre
basis used in R4's referee report, section 1, step 9.4, one has the
signed formula
\[
 T_\sigma=\kappa(\sigma)\rho_{E_3,12}(\sigma),\qquad
 \kappa(\sigma)=\sigma(r_0)/r_0\in\{\pm1\},
 \quad \sigma\in G_{\mathbf Q^{\rm ab}}.
\]
It follows by writing the equianharmonic Legendre fiber as the
quadratic twist of $y^2=x^3+1$ by $1/\sqrt{-3}$; the square root
has the same character as $r_0$, since $\sqrt2$ and $\zeta_8$
belong to $\mathbf Q^{\rm ab}$. Reduction modulo three and
quotienting by $\langle-I\rangle$ removes exactly this quadratic
twist. The resulting projective torsion representation is therefore
the projection of $T$ to $C_6/\langle-I\rangle$, with no further
unspecified normalization. Its image has order three, so the
Galois frame image has order divisible by three.

Let $r_0=(3/4)^{1/4}$ and $M=\mathbf Q(i,r_0)$, the field
identified in the refereed proof of R4. The polynomial
$X^4-12$ of $2r_0$ is Eisenstein at three, and
$\mathbf Q(r_0)$ is real. Therefore $[M:\mathbf Q]=8$, with
the dihedral action $r_0\mapsto i r_0$, $i\mapsto i$ and
complex conjugation. Its maximal abelian subfield is
$\mathbf Q(i,r_0^2)=\mathbf Q(i,\sqrt3)$. Hence an element of
$G_{\mathbf Q^{\rm ab}}$ sends $r_0$ to $-r_0$. The verified
sign convention of R4 makes both parities one for that element.
Its $W$ component is an odd power of $S$ and has order four.
The Galois image in the cyclic group $J$ thus has order divisible
by both three and four, so it is all of $J$.

For the element of R2, the corrected coordinates are $(-I,I)$.
It is outside $J$. Consequently the image of $\mathcal H$
contains the order-12 subgroup $J$ and its other coset, and is
all of $C_6\times C_4$. The group enumeration is independently
checked in \texttt{check\_cm\_replacement.py}. \hfill$\square$

\begin{statement}{B.15}{Proved}{the arithmetic replacement lifts}
\citep[Section 8.5]{LiftHainMatsumoto2003}
\end{statement}
Let $N=\ker q_{12}\subset\widehat B$. There is a continuous
positive peripheral automorphism of
$\widehat B/R_{\rm nil}(N)$ with the LM7 D.46 component
description and specified conjugators. It lifts by R2 to an
element of $\mathcal H$ with projective pair $(I,[i])$, parities
$(0,1)$ and $\Delta=1$. We can prove the needed arithmetic
input without assuming LM7 D.35's entire 2-adic image theorem.

\emph{Proof.}
Put $k=\mathbf Q(\zeta_3)$ and
$a=1-2\zeta_3^2=3+2\zeta_3=2+\sqrt{-3}$. Counting
$E_3(\mathbf F_7)$ gives 12 points, so Frobenius at either
prime of $k$ over seven has characteristic polynomial
$X^2+4X+7$. The automorphism $(x,y)\mapsto(\zeta_3x,y)$
makes $T_2E_3$ a free rank-one module over
$\mathbf Z_2[\zeta_3]$. Frobenius commutes with this module
structure and is multiplication by one of $-a,-\bar a$.
Choose the prime, or conjugate it, so that it is $-a$.

All two-torsion of $E_3$ is $k$-rational. Its full 2-division
extension over $k$ is therefore pro-two: its image is in the
pro-two kernel of reduction modulo two. Its intersection with
$k(\sqrt[3]4)=k(E_3[3])$ is $k$, since the latter extension is
cyclic of degree three. The resulting product surjection of
Galois groups permits an actual $g\in G_k$ with 2-adic image
$-a$ and trivial action on $E_3[3]$. This uses only one actual
Frobenius and the coprime disjointness of these extensions.

Choose $\tau\in I_3$ inducing the nontrivial automorphism of
$k$, which is ramified at three, and put $h=[g,\tau]$.
The full cyclotomic character of a commutator is one. The
normalized action of $I_3$ on the free pro-two sphere group is
trivial by good reduction, with the integral tangential base
point $\overrightarrow{01}$. This exact automorphism statement
is stated in Hain--Matsumoto, section 8.5 immediately before
Theorem 8.7, for $N=1,\ell=2$. Consequently the normalized
action of $h$ on the entire pure pro-two group is the identity.
On $T_2E_3$, conjugation by $\tau$ conjugates the CM coefficient
ring. Thus
\[
 \rho_{E_3,2}(h)=a/\bar a=(5+8\zeta_3)/7
                    =(-3-8\zeta_3^2)/7.
\]
It is $-I$ modulo four. Also $h$ acts trivially on $E_3[3]$,
since $g$ does. The signed frame convention of R4, which is
$I$ modulo four, therefore gives $T_h=-I$ modulo three.
Its first parity is one and, by R4, its second parity is one.

We check that the induced action on $Q$ is literally the identity.
For every multiplier-one pair, the pure $r$-frame $A$ belongs
to $F'$. Indeed $f=rAr^{-1}A^{-1}$, so
$(r-1)A^{\rm ab}=0$; the action of $r-1$ on
$F^{\rm ab}=\widehat{\mathbf Z}^2$ is injective, with determinant
three. Hence $q_3(A)$ is in the quaternion derived subgroup of
$\mathrm{SL}_2(\mathbf F_3)$. In its abelianization $C_3$, the
identity $T=q_3(A)^{-1}U^{8\rho}$ gives $[T]=2\rho[U]$.
Since $T_h=-I$ has trivial class there, $\rho_h\equiv0\pmod3$.
The shear $U^{8\rho_h}$ is consequently $I$ modulo twelve.
R3 now gives $q_{12}\Phi_h=q_{12}$. Its frame coordinates in
$Q$ are $a_Q=[5I]$ and some $b_Q$ with inverse reducing to
$\pm S$ at three and to $I$ at four.

For each prime $p$, put $E_p=\widehat B/O^p(N)$. The actual
fiber product theorem of R2 gives
\[
 \widehat B/R_{\rm nil}(N)=\prod_Q E_p.
\]
Because $\Phi_h$ preserves $N$ and is the identity on $Q$,
the tuple with $\Phi_h$ in the two-component and identity in
all odd components is a continuous automorphism of this
fiber product fixing $\tau$. We must specify its peripheral
conjugators; a product of automorphisms alone would not suffice.

In $E_2\to Q$, the element $a_Q^{-1}$ centralizes $r_Q$.
Coprime complement conjugacy for the order-three action lifts
it to $c_2\in C_{E_2}(r)$. Therefore $A_{h,2}c_2$ is a
conjugator of $r$ with $Q$-coordinate one. At every odd $p$,
the element $b_Q$ centralizes $s_Q$ and lifts by the same
argument to $b_p\in C_{E_p}(s)$. The two actual tuples
\[
 A^\dagger=(A_{h,2}c_2,(1)_{p\ne2}),\qquad
 \beta^\dagger=(\beta_{h,2},(b_p)_{p\ne2})
\]
have common $Q$-coordinates $1,b_Q$. They belong to the image
of $F$ and are the required positive peripheral conjugators;
the cusp conjugator is one. Thus every hypothesis of R2 is
verified, including the prime-two and prime-three cases and
the mixed finite base $Q$.

R2 lifts the automorphism with these conjugators to $\widehat B$
and keeps $\tau$ fixed. The resulting parameter $f=s\Phi(s)$
is trivial in every absolute pure pro-$p$ quotient: for $p=2$
this was proved for $h$, and for odd $p$ the chosen component
is the identity. More explicitly, an identity pure action gives
$f_p\in C_{F^{(p)}}(y)=\langle y\rangle$, while the actual
torsion relation $(sf_p)^2=1$ makes its abelian coordinate
$(a,-a)$. An element of $\langle y\rangle$ has zero first
coordinate, so $a=0$ and this element is one.
In particular $f\in F'$. The torsion word
calculation in P.9 now proves both cycles, and the restriction
is an actual automorphism. Hence it is an element of $\mathcal H$.
Its induced action on $Q$ is still literally the identity.
By R3 its shear $U^{8\rho}$ therefore centralizes the image of
$F$; modulo three that image is $\mathrm{SL}_2(\mathbf F_3)$.
The intersection of its center with the upper unipotent group
is trivial, so $\rho\equiv0\pmod3$ and the shear is $I$
modulo twelve. This distinguishes the corrected shear from the
separately specified cusp conjugator one.
Its retained $Q$-frames have $A_Q=1$, $\beta_Q=b_Q$, and
corrected $V_Q=1$, giving $T=I$, $W=\pm S$. Projectively this is the
pair $(I,[i])$ claimed in LM7 D.46, and its parities are $(0,1)$.

In the graph basis in which $\zeta_3^2$ acts by
$\left(\begin{smallmatrix}0&1\\-1&-1\end{smallmatrix}\right)$,
the computed scalar is
$\left(\begin{smallmatrix}11&8\\8&3\end{smallmatrix}\right)$
modulo sixteen. This and the Frobenius point count are recorded
by \texttt{check\_cm\_replacement.py}. The construction does
not assert that this element satisfies the pentagon.
\hfill$\square$

\begin{statement}{B.16}{Proved}{the discrepancy is independent of absolute pronilpotent data}
\end{statement}
Let $v:\widehat{\mathrm{GT}}_0\to V$ be the product of its
specified absolute pro-prime, pronilpotent and associated graded
actions, with $V$ its image. Then
\[
 (v,\Delta):\widehat{\mathrm{GT}}_0
                  \longrightarrow V\times\mathbf F_2
\]
is a continuous epimorphism. The Galois image has zero second
coordinate. With $V_{\mathcal H}=v(\mathcal H)$, the stronger
map $(v,\epsilon_3T,\epsilon_2W):\mathcal H\to
V_{\mathcal H}\times\mathbf F_2^2$ is onto.

\emph{Proof.}
If $U$ is any subgroup of $F$ and $q:F\to D$ has finite
nilpotent image, then $q(U)$ is nilpotent. Hence
$R_{\rm nil}(U)\subseteq R_{\rm nil}(F)$. The R2 witness,
which is identity on $\widehat B/R_{\rm nil}(F'F^3)$, is
therefore identity on every absolute pro-$p$ quotient of $F$,
on their product, and on all the associated graded quotients.
It lies in $\ker v$ and has $\Delta=1$. Multiplying any chosen
lift of a point of $V$ by this element toggles the second
coordinate. This proves the first epimorphism.

The construction in B.15 is also identity on every absolute
pure pro-$p$ quotient and has parity pair $(0,1)$; the R2
witness has pair $(1,0)$. Both lie in $\ker v$. Their products
realize every parity pair in that kernel, proving the stronger
assertion on $\mathcal H$. The restriction on the Galois image
is exactly R4. \hfill$\square$

The word ``quotients'' in this assertion refers to the geometric
actions induced from the absolute pro-prime and pronilpotent
quotients of $F$, and their associated graded objects. If it
instead means all pro-prime quotients of the profinite group
$\widehat{\mathrm{GT}}_0$ itself, the proposed assertion is false:
the nonzero character $\Delta$ is already a quotient of that group
of order two, hence factors through its maximal pro-two quotient.
In that interpretation its value is determined by $V$, so the
image is a graph rather than $V\times\mathbf F_2$.

\begin{statement}{B.17}{Refuted}{the signs of LM8 B.31 are not the CM frame parities}
\end{statement}
The proposed identification in B.6(e) is false. The linked and
independently signed actual sources of LM8 B.31 coincide for
an elementary abelianization reason which concerns the powers
of torsion generators, whereas $\Delta$ concerns their
conjugating frames.

\emph{Proof and replacement.}
Let $C=\widehat{\mathrm{SL}_2(\mathbf Z)}$ with
$\rho^6=\sigma^4=1$, $\rho^3=\sigma^2$, and $e_0=\sigma\rho$.
Its abelianization is $\mathbf Z/12$, with $[\rho]=2$,
$[\sigma]=-3$, $[e_0]=-1$. An actual automorphism with
$u(e_0)=e_0^\mu$, $u(\rho)\sim\rho^{e_r}$ and
$u(\sigma)\sim\sigma^{e_s}$ therefore satisfies
\[
              e_r\equiv\mu\pmod6,\qquad
              e_s\equiv\mu\pmod4.
\]
The two exponent signs are forced by $\mu$; at multiplier one
they are both positive. The same calculation on the projective
modular group forces the corresponding positive exponent
classes there. Nevertheless the R2 witness has multiplier one
and CM frame parities $(1,0)$. Thus the two kinds of signs
cannot be identified. The strongest true statement is that
the exponent signs describe the cyclotomic power action,
while the CM parities are residual conjugator information
after that power action and the cusp shear have been fixed.
This addresses the proposed comparison with section 49.IV--V,
cote 140-3, scans 530, 549 and 559--560 as cited in LM8 B.31;
the original scans are not supplied here, so it makes no new
palaeographic assertion about them. \hfill$\square$

\subsection{Metabelian units and arithmetic compatibility}


\begin{statement}{B.18}{Known}{the precise metabelian input}
\citep{LiftChenDeligne2017}
\end{statement}
Put $M=F/F''$, $A=M^{\rm ab}=\widehat{\mathbf Z}^2$,
$R=\widehat{\mathbf Z}[[A]]$, and $c=[x,y]$. The continuous
Alexander module is $M'=Rc$, free of rank one. Moreover
\[
 \operatorname{IOut}(M)=R^{\times,\epsilon=1}/A,
 \qquad
 \operatorname{Out}(M)=
 (R^{\times,\epsilon=1}/A)\rtimes\operatorname{GL}(A),
\]
with the natural action on the completed group ring. The
determinant of an IA automorphism is its multiplier on $Rc$;
an IA automorphism is inner exactly when this determinant is a
monomial in $A$. These are Chen--Deligne, Corollary 4.1.5,
Proposition 4.3.7 and Theorem 5.2.4. Their Corollary 4.1.6 also
gives $Z(M)=1$. The cited statements concern this actual
profinite quotient, rather than the completion of a discrete
formula with an unverified comparison map.

\begin{statement}{B.19}{Proved}{LM6 A.10 and A.49 are the units of LM8 A.22}
\end{statement}
There is a continuous injective homomorphism with closed image
\[
 \widehat{\mathbf Z}^{\times}\longrightarrow
 C_{\operatorname{Out}_{\rm lac}(M)}(G_K),\quad
 t\longmapsto[\phi_t]
\]
for every number field $K$. Its image is IA, centralizes the
geometric $S_3$, and is identity in every maximal pro-prime
quotient. The coefficient-five element $t_5=-1$, $t_p=1$ for
$p\ne5$, is exactly the involution of LM8 A.22. The analogous
choices for arbitrary sets of primes $p\ge5$ give LM6 A.10.

\emph{Proof.}
For $p\ne3$, let $A^{(p')}\to C_3$ send both named generators
to the same generator $t_0$, and let $H$ be its kernel.
Normalized averaging on every finite image of $H$ gives a
compatible idempotent $e_H\in\mathbf Z_p[[A^{(p')}]]$, because
all its averaging denominators are prime to $p$. Define
\[
 e_p=e_H\left(1-\frac{1+t_0+t_0^2}{3}\right)
\]
in the coefficient-$p$ factor, and zero in all other factors.
It is nonzero and has augmentation zero. On its summand,
the prime-to-$p$ parts of $X,Y,XY$ are nontrivial cubic roots;
therefore $X-1,Y-1,XY-1$ are units there. The kernel $H$ and
the pair of nontrivial characters are invariant under scalar
units and under the peripheral $S_3$. This construction applies
also at $p=2$.

For $p=3$ use $B=(\mathbf Z/4)^2$ and the six characters
\[
 \mathcal S=\{(a,b):a\ne0,\ b\ne0,\ a+b\ne0\pmod4\}.
\]
Their Fourier projector belongs to $\mathbf Z_3[B]$: conjugate
characters pair, and its coefficients have denominator dividing
eight. Average over the kernel of $A^{(3')}\to B$ to lift it
to an idempotent $e_3$ in the coefficient-three factor of $R$.
The set $\mathcal S$ is stable under $(a,b)\mapsto(b,a)$,
$(a,b)\mapsto(-a-b,b)$ and $(a,b)\mapsto(-a,-b)$; hence $e_3$
has the same invariance properties. Its augmentation is zero,
and $X-1,Y-1,XY-1$ are units on its summand.

For $t=(t_p)_p$ define, component by component,
\[
 u_t=1+\sum_p(t_p-1)e_p,\qquad
 r_t=\frac{u_t-1}{(X-1)(Y-1)},\qquad
 s_t=\frac{X(Y-1)r_t}{XY-1}.
\]
Each fraction is computed only on the indicated idempotent
summands and is extended by zero elsewhere. Put
$f_t=c^{r_t}$, $h_t=c^{s_t}$, where the exponents denote the
continuous $R$-module structure. The actual endomorphism
\[
 \phi_t(x)=x,\qquad \phi_t(y)=f_t^{-1}yf_t
\]
exists by the universal property of the free profinite
metabelian group. Direct multiplication gives
\[
 \phi_t(c)=c^{u_t},\qquad
 \phi_t(xy)=c^{X(Y-1)r_t}xy=h_t^{-1}xyh_t.
\]
It preserves all three labelled peripheral classes. Since it
is IA, it acts $R$-linearly on $M'$, and the conjugator for
$\phi_t\phi_v(y)$ has exponent
\[
 r_t+u_tr_v
   =\frac{u_tu_v-1}{(X-1)(Y-1)}=r_{tv}.
\]
Consequently $\phi_t\phi_v=\phi_{tv}$ and $\phi_{t^{-1}}$
is its inverse; these are actual automorphisms. In particular
$t_p=-1$ gives determinant $1-2e_p$ and an involution.

Arithmetic acts on $A$ as the scalar $\chi(\sigma)$. In B.18,
conjugation by an arbitrary pair $(v,g)$ sends $(u,1)$ to
$(g(u),1)$, because the IA unit group is abelian. Invariance
of every $e_p$ thus proves outer commutation with $G_K$ and
with $S_3$. This step does not assert equality of arbitrarily
chosen representatives in $\operatorname{Aut}(M)$.

Projection to the coefficient-$p$ maximal pro-$p$ Alexander
module kills every $e_q$: for $q\ne p$ its coefficient is
killed, and for $q=p$ its nontrivial prime-to-$p$ character
becomes trivial. Thus the actual automorphism becomes the
identity in each maximal pro-prime quotient. If $u_t$ were a
monomial, these projections would make each pro-$p$ component
of that monomial trivial; the monomial would be one. The
nonzero torsion-free module $e_pR$ would then force $t_p=1$
for all $p$. B.18 therefore proves outer injectivity. Every
finite group quotient uses finitely many coefficient primes
and residue digits, giving continuity. Compactness gives a
closed image.

The new independent program \texttt{check\_metabelian\_units.py}
checks the coefficient-three Fourier idempotent, its square,
augmentation and all stated symmetries exactly. It also constructs
the finite quotients of orders 48, 75 and 324 and checks generation,
every multiplication pair, all possible inner conjugators and the
three peripheral identities. These checks include both primes
excluded from the elementary cubic-character description at five.
They certify the displayed finite images; the inverse-limit
construction above proves the infinite assertion. \hfill$\square$

\begin{statement}{B.20}{Proved}{the center-free comparison needed one derived step higher}
\end{statement}
For every $m\ge2$, the group $S_m=F/F^{(m)}$ is center-free.
Hence its arithmetic extension over a number field $K$ is
canonically the pullback
\[
 E_m\simeq\operatorname{Aut}(S_m)
             \times_{\operatorname{Out}(S_m)}G_K.
\]
An outer automorphism centralizing $G_K$ therefore induces an
actual continuous automorphism of $E_m$ over $G_K$.

\emph{Proof.}
The base case is B.18. Inductively a central element of $S_m$
lies in $F^{(m-1)}/F^{(m)}$. Let $H$ range over the open normal
subgroups of $F$ containing $F^{(m-1)}$. Its image $z_H$ in
$H^{\rm ab}$ is fixed by $F/H$ and maps to zero in $F^{\rm ab}$.
The transfer identity therefore gives
\[
 [F:H]z_H=\operatorname{Ver}_{F/H}(0)=0.
\]
Profinite Nielsen--Schreier makes $H^{\rm ab}$ a finite-rank
free profinite abelian group, which has no nonzero integer
torsion; hence $z_H=0$.
Moreover $\bigcap_HH'=F^{(m)}$: an element of $F^{(m-1)}$
outside $F^{(m)}$ survives in a finite quotient of $S_m$;
the inverse image of the abelian image of $F^{(m-1)}$ is one
such $H$, and its derived subgroup cannot contain the chosen
element. Thus the central element is trivial.

The map from $E_m$ to the displayed pullback sends an arithmetic
element to its conjugation action and its image in $G_K$.
Its kernel is $Z(S_m)=1$, and its image contains every fiber:
any two representatives of a prescribed outer action differ by
an inner automorphism. It is therefore a continuous isomorphism
of profinite groups. Conjugating the first component by a
representative of the commuting outer class defines the claimed
extension automorphism. \hfill$\square$

\begin{statement}{B.21}{Proved}{audit of LM6 A.56, including the constant twists}
\citep{LiftSerre2007,LiftMilneCFT2020,LiftMilneLEC2013}
\end{statement}
For every number field $K$, the image in $\operatorname{Out}(M)$
of the centralizer of $G_K$ in $\operatorname{Out}(S_3)$ contains
no nontrivial IA class. Here $S_3=F/F'''$ denotes the third
derived quotient, not the symmetric group. In particular every
nonidentity unit in B.19 is obstructed one derived step higher
if commutation with $G_K$ is required.

\emph{Proof.}
Let $[\psi]$ be such a class and let $u\in R^{\times,\epsilon=1}$
represent its metabelian determinant. Write $I=\ker\epsilon$.
The continuous map $u\mapsto u-1\bmod\overline{I^2}$ identifies
$I/\overline{I^2}$ with $A$: $X^aY^b$ has linear term
$a(X-1)+b(Y-1)$, and the square-zero ring
$\widehat{\mathbf Z}\oplus A$ proves the converse independence.
Multiplying by the inverse monomial with this linear term
normalizes $u$ to a unit $u_0$ with zero linear term.

Outer commutation says $u_0/\sigma_R(u_0)\in A$ for every
$\sigma\in G_K$. This monomial has zero linear term and is
therefore one. Thus $\sigma_R(u_0)=u_0$. Inner adjustments lift
through $S_3\to M$, so choose an actual representative whose
metabelian determinant is $u_0$. By B.20 it extends to an
automorphism $\beta$ of $E_3$ over $G_K$.

This choice makes $\beta$ the identity on $E_1$. Indeed if
$b_\sigma$ is conjugation by an arithmetic lift on $M$, its
action on $Rc$ is $r\mapsto\sigma_R(r)v_\sigma$. The commutator
of our representative with $b_\sigma$ acts on $Rc$ by
$u_0/\sigma_R(u_0)=1$. It is inner, say by $d_\sigma\in M$.
An inner action on $Rc$ multiplies by the monomial
$d_\sigma^{\rm ab}$, and its faithfulness gives $d_\sigma\in M'$.
The correction of each arithmetic lift thus vanishes in $E_1$;
the geometric abelian action was already the identity.

Suppose $u_0\notin A$. The closed procyclic subgroups
$C=\widehat{\mathbf Z}c$ and $D=\widehat{\mathbf Z}u_0c$ in
$M'$ cannot be conjugate. A conjugator acts by a monomial
$a\in A$, so conjugacy would give $u_0=\lambda a$ with
$\lambda\in\widehat{\mathbf Z}^{\times}$; augmentation forces
$\lambda=1$. Compactness detects their nonconjugacy in a
finite quotient: otherwise the nonempty closed sets of
conjugators modulo every open normal subgroup would have
the finite-intersection property and yield an actual conjugator.
Intersecting the finite automorphism-orbit of a detecting
open normal subgroup produces a stable finite metabelian
quotient $q:M\to G$ with induced automorphism $\gamma$ and
nonconjugate equal-order cyclic groups
$J=\langle q(c)\rangle$, $J'=\gamma(J)$. For $g=q(c)$,
\[
             \#(G/J)^g>0,\qquad \#(G/J')^g=0.
\]

After a finite extension of $K$, the geometric $G$-torsor
$Y\to U_0=\mathbf P^1\setminus\{0,1,\infty\}$ descends with
constant deck group and an arithmetic quotient $G\times G_K$.
On this quotient $\beta$ has form
$(h,\sigma)\mapsto(\gamma(h)z_\sigma,\sigma)$, with
$z:G_K\to Z(G)$ a continuous homomorphism. A further finite
extension kills $z$. We may also replace $K$ by its Galois
closure; this does not change the identity on $E_1$.

Hilbert irreducibility supplies $a\in U_0(K)$ whose fiber
$Y_a$ is a $G$-Galois field extension $L/K$; this is precisely
Serre, Proposition 3.3.1 and Theorem 3.4.1, in the opened
second edition. Lift $g$ to the normal closure of $L/\mathbf Q$.
Chebotarev, as stated in Milne, \emph{Class Field Theory},
Chapter V, Theorem 3.23, gives a rational prime $q$ split in
$K$ and a degree-one place $v$ whose fiber Frobenius is $g$.
Exclude the finitely many primes of bad marked reduction,
those dividing $|G|$, and those where $a$ or $1-a$ is not a
unit. Put $N=q-1$ and $K'=K(\mu_N)$. The place extends with
residue field $\mathbf F_q$, since $q\equiv1\pmod N$.

Use the actual constant-twisted abelian cover
\[
 B:\quad z_1^N=t/a,\qquad z_2^N=(1-t)/(1-a).
\]
Its arithmetic monodromy factors through $E_1$, so $\beta$
preserves it exactly, including its two Kummer constants.
Its only nonempty rational fiber over $\mathbf F_q$ is above
$a_0=a\bmod v$, where it has $N^2$ points. Therefore the
covers $Z=(Y/J)\times_{U_0}B$ and
$Z'=(Y/J')\times_{U_0}B$ satisfy
\[
 \#Z(\mathbf F_q)=N^2\#(G/J)^g>0,
 \qquad \#Z'(\mathbf F_q)=N^2\#(G/J')^g=0.
\]

On the other hand their geometric monodromy is metabelian.
Every component subgroup $H\subset F$ contains $F''$, hence
$F'''\subset H'$. Thus the quotient $F/F'''$ retains exactly
the full abelianization $H^{\rm ab}$ required for its first
cohomology. The automorphism $\beta$ identifies the arithmetic
cover data, the component permutation modules, and these
abelianizations, giving equivariant identifications of both
$H^0$ and $H^1$ of $Z_{\overline K}$ and $Z'_{\overline K}$.
For any coefficient prime different from $q$, good marked
reduction transports these representations to the special
fiber. The normalized products have good marked reduction:
locally at a boundary point their equations are tame Kummer
equations, and normalization replaces the indices by their
least common multiple, still prime to $q$. Poincare duality
gives equality on $H_c^1$; $H_c^0=0$ and $H_c^2$ is the
common component module twisted by $-1$. The trace formula
therefore forces equal point counts, a contradiction. The
cohomology facts used here are Milne, \emph{Lectures on Etale
Cohomology}, Theorems 24.1 and 29.4, together with proper
smooth base change for the smooth compactifications and their
finite etale boundaries.

It follows that $u_0\in A$, and B.18 makes the metabelian IA
outer class trivial. Notice the exact scope: the argument
does not prove that the kernel of
$C_{\operatorname{Out}(F/F''')}(G_K)\to\operatorname{Out}(M)$
is trivial. \hfill$\square$

\begin{statement}{B.22}{Proved}{audit of LM6 A.57--A.58}
\end{statement}
The abelian image of the peripheral metabelian centralizer
of $G_K$ is exactly the geometric symmetric group. Consequently
\[
 \operatorname{image}\left(
 C_{\operatorname{Out}_{\rm lac}(F/F''')}(G_K)
       \longrightarrow\operatorname{Out}(M)\right)
       =S_3^{\rm geom},
\]
and the same holds for every higher derived length projecting
to $M$.

\emph{Proof.}
An abelian automorphism preserving the three primitive
peripheral lines is a scalar unit times a geometric permutation:
fixing the first two lines makes it diagonal and fixing their
sum makes the diagonal entries equal. Divide by the geometric
permutation and write its common multiplier $\rho$.

We need an exact arithmetic normalization. In $M$ the normalizer
of $\langle x\rangle$ is $\langle x\rangle$. Indeed a centralizer
element has zero $y$-coordinate by the free class-two quotient,
and hence is $x^a c^r$. Commutation says $(X-1)r=0$. In every
coefficient-$p$ factor, refining the first cyclic group coordinate
by $p^j$ makes an $X$-invariant coefficient project to $p^j$
equal summands; as $j$ is arbitrary it is zero. This proves
injectivity of $X-1$ on the full completed ring. Normalizing
$\langle x\rangle$ already forces exponent one by abelianization,
so the normalizer equals the centralizer. The same holds for $y$.

The tangential sections at $0$ towards $1$ and at $1$ towards
$0$ have the same image in $E_1=A\rtimes G_K$: on the covers
$t^{1/n}$ and $(1-t)^{1/n}$ their leading coefficients or values
are one, and choose the compatible roots one. Their arithmetic
normalizers are respectively $\langle x\rangle\rtimes G_K$
and $\langle y\rangle\rtimes G_K$. Write an extension
automorphism obtained from B.20 as
$(v,g)\mapsto(\rho v+b(g),g)$ on $E_1$. Peripheral preservation
and the two normalizer equalities imply
$b_2(g)=(1-\chi(g))(a_0)_2$ and
$b_1(g)=(1-\chi(g))(a_1)_1$, where $a_i$ are the abelian
coordinates of its two geometric cusp conjugators. Thus $b$
is an actual coboundary. An inner adjustment gives exactly
$(v,g)\mapsto(\rho v,g)$.

Let $q$ be an odd rational prime split in a Galois closure of
$K$, $N=q-1$, and pass to $K'=K(\mu_N)$ at a degree-one place
over $q$. For each integer $a\in\{2,\ldots,q-1\}$ use
\[
 B_a:\ z_1^N=t/a,\ z_2^N=(1-t)/(1-a).
\]
If $r=\rho\bmod N$, the exact abelian action sends its
arithmetic cover subgroup to that of
\[
 B'_a:\ w_1^N=t/a^r,\ w_2^N=(1-t)/(1-a)^r.
\]
This follows by applying $(v,g)\mapsto(\rho v,g)$ to the
equations $v=(\kappa_N(a)(g),\kappa_N(1-a)(g))$ defining the
cover subgroup. Their geometric monodromy is abelian, so
$F''$ lies in each component subgroup's commutator and already
$E_2$ retains its full first cohomology. The same cohomology
and trace argument as B.21 forces equal point counts. Both
covers have good marked reduction, being Fermat curves of
degree $q-1$ with unit coefficients and their finite etale
boundaries removed. The first has $N^2$ rational points;
the second has $N^2$ precisely when
\[
                       a^r+(1-a)^r=1\quad\hbox{in }\mathbf F_q.
\]
This identity therefore holds for all $a\ne0,1$. The map
$f(0)=0$, $f(x)=x^r$ on $\mathbf F_q^\times$ is multiplicative
and bijective. Taking $a=x/(x+y)$ proves additivity for nonzero
$x,y,x+y$. The remaining cases follow from $f(-1)=-1$, since
$r$ is odd. Thus $f$ is an automorphism of the prime field,
and is the identity. Hence $\rho\equiv1\pmod{q-1}$.
For every $n$, Chebotarev permits such a $q$ also split in
$\mathbf Q(\mu_n)$, so $\rho\equiv1\pmod n$; thus $\rho=1$.

The projected centralizer at length three consequently has
the abelian action of one geometric permutation. Dividing by
it leaves an IA class, trivial by B.21. The geometric six
permutations are defined over $\mathbf Q$ and lift at every
derived level. This proves equality, including higher levels
by projection to length three. \hfill$\square$

\begin{statement}{B.23}{Proved}{reconciliation with the three weight-five directions}
\end{statement}
B.11 concerns arbitrary multiplier-one peripheral automorphisms
of the full free profinite group; it imposes no commutation
with arithmetic. B.19 concerns actual metabelian outer classes
which commute with arithmetic, and B.21 excludes lifting these
classes while retaining that extra condition. Their determinants
lie on nontrivial prime-to-$p$ character summands and have zero
image in every absolute pro-$p$ quotient, so they are not any
of the nonzero weight-five directions. Thus there is neither
a conflict of categories nor a counterexample to R1. R1 retains
three peripheral conjugators, but its hypotheses do not include
Galois commutation. At primes two and three B.19 and B.21 apply
unchanged; the invariant pro-$\ell$ lifting step of P.8 is
only asserted for $\ell\ge5$, exactly where division by six
and the two coprime complement arguments are valid.

\subsection{The manuscript dictionary}


\begin{statement}{B.24}{Proved}{the numerical beta is half-cyclotomic}
\end{statement}
Use the exact conventions of cote 145, scans 7--12:
$x=l_0$, $y=l_1$, $z=(xy)^{-1}$,
$\rho(x)=y$, $\rho(y)=z$, and
$\sigma_\infty(x)=y$, $\sigma_\infty(y)=x$.
Let $u(x)=x^\lambda$ and
\[
 u\rho=\operatorname{Inn}(g)\rho u,
 \qquad
 \sigma_\infty(g^{-1})=x^\beta g y^\beta.
\]
The last relation is written without a disputed reading in the
supplied transcription at scan 12. If $g^{\rm ab}=(a,b)$,
then
\[
 \beta=-(a+b)=\frac{\lambda-1}{2}.
\]
In particular this numerical $\beta$ is zero on $\mathcal H$.
It is not the $s$-frame coordinate $\rho_F$ of R3.

\emph{Proof.}
Abelianizing the scan-12 equation gives
$(-b,-a)=(a+\beta,b+\beta)$, so $\beta=-(a+b)$.
To identify the sum, work in the actual continuous free
class-two quotient of $F$, writing $c=[x,y]$ and normal forms
$x^ay^bc^d$. Its multiplication law is
\[
 (a,b,d)(a',b',d')=(a+a',b+b',d+d'-a'b).
\]
These identities hold with profinite coefficients by continuity.
Since $u(y)=gy^\lambda g^{-1}$, one obtains
\[
 u(z)=x^{-\lambda}y^{-\lambda}
                        c^{-a\lambda-\lambda^2}.
\]
On the other hand, the equation for $u\rho$ evaluated at $y$
gives $u(z)=g\rho(g)z^\lambda\rho(g)^{-1}g^{-1}$.
Here $g\rho(g)$ has abelian coordinate $(a-b,a)$, and
\[
 z^\lambda=x^{-\lambda}y^{-\lambda}
                         c^{-\lambda(\lambda+1)/2}.
\]
The second expression for $u(z)$ therefore has central
coordinate $-\lambda(\lambda+1)/2+b\lambda$.
Equality gives
$-\lambda(a+b)=\lambda(\lambda-1)/2$; cancel the unit
$\lambda$ to obtain the claimed formula. Since $\lambda$
is a profinite unit, $\lambda-1$ is divisible by two in
$\widehat{\mathbf Z}$, and division by two is unique.
The symbolic class-two calculation and the finite
$\mathrm{SL}_2(\mathbf F_3)$ check below are recorded in
\texttt{check\_manuscript\_cocycle.py}.

For a concrete separation from R3, take a Galois element in
$G_{\mathbf Q^{\rm ab}}$ with nontrivial action on
$\mathbf Q(\zeta_3,\sqrt[3]4)/\mathbf Q(\zeta_3)$; its
existence was proved in B.14. The exact signed frame comparison
makes $T$ nontrivial in
$\mathrm{SL}_2(\mathbf F_3)^{\rm ab}=C_3$. The pure $r$-frame
belongs to $F'$ by the argument in B.15, so in this abelianization
\[
                   [T]=[U^{8\rho_F}]=2\rho_F[U].
\]
Here $[U]$ generates $C_3$. Thus $\rho_F\not\equiv0\pmod3$,
whereas the manuscript's $\beta=0$, since $\lambda=1$.
The finite matrix check verifies that the derived subgroup
has order eight and that both $U$ and the equianharmonic
order-three generator have nontrivial classes. This refutes
$\beta=\rho_F$, including on the multiplier-one subgroup.
\hfill$\square$

\begin{statement}{B.25}{Proved}{the corrected injective pair dictionary}
\end{statement}
The manuscript's gauge class of $g$, together with $\lambda$,
determines the normalized GT pair injectively. Put
$m=(\lambda-1)/2$. In the gauge where $a=0$, the exact formula is
\[
                   g=f^{-1}y^{-m},\qquad f\in F'.
\]
For an arbitrary representative with $g^{\rm ab}=(a,b)$, the
gauge-invariant formula is
\[
                          f=y^b g^{-1}x^a.
\]
Consequently the formula $g=f^{-1}$ is correct on
$\mathcal H$ after normalization, but needs the explicit
half-cyclotomic factor at general multiplier. This corrects
the unrestricted reading of the dictionary in R11.

\emph{Proof.}
Replacing $u$ by $\operatorname{Inn}(x^\alpha)u$ changes $g$
to $x^\alpha gy^{-\alpha}$ and its abelian coordinate to
$(a+\alpha,b-\alpha)$. There is exactly one gauge making
the first coordinate zero, namely $\alpha=-a$. By B.24 the
second coordinate is then $-m$. Define $f=y^{-m}g^{-1}$
in that gauge; it lies in $F'$, and
$u(y)=f^{-1}y^\lambda f$, while $u(x)=x^\lambda$.
Returning to the original representative gives exactly
$f=y^b g^{-1}x^a$. Substituting the gauge transformation
in this expression shows it is unchanged, with the order
of every factor retained.

The normalized pair determines the images of the two free
generators and therefore the continuous automorphism and its
outer class. Conversely the normalization just constructed
recovers that pair from the gauge class and multiplier, proving
the stated injectivity on the manuscript's $\Gamma$, identified
with $\widehat{\mathrm{GT}}_0$ under the marked conventions.
The dictionary is a conversion of coordinates, not a claim
that the set of gauge classes itself carries the componentwise
group law. For complex conjugation, $\lambda=-1$, $m=-1$,
and the normalized modern $f=1$ gives $g=y$, exactly the
manuscript's scan-9 and scan-11 value and scan-12 $\beta=-1$.
The gauge preserves $a+b$, so it cannot remove this extra
$y$ factor. \hfill$\square$

\subsection{Sewing and comparison source checks}


\begin{statement}{B.26}{Known}{the selected sewing stabilizer}
\citep{LitBoggi2026v6,LitMargalitWinarski}
\end{statement}
Let $S$ be the closed oriented surface of genus two and let $\gamma$ be
a fixed oriented nonseparating simple closed curve. All mapping classes
preserve the orientation of the surface. The subscript $\vec\gamma$
requires preservation of the orientation of the curve as well.
Let $t$ be its positive Dehn twist. For the compact genus-one surface
$S_1^2$ with its two boundary components fixed pointwise, write $d$ for
the first positive boundary twist. In the standard positive three-chain
presentation, there is an identification
\[
 \Gamma(S_1^2,\partial S_1^2)=B_4\times\langle d\rangle,
 \qquad \delta_1=d,\qquad \delta_2=z_4d^{-1},
 \qquad z_4=(\sigma_1\sigma_2\sigma_3)^4.
\]
Margalit--Winarski, Section 6, gives $\operatorname{Mod}(S_1^2)
\cong B_4\times\mathbf Z$, with the extra generator a boundary twist;
the positive three-chain relation identifies the second boundary twist
as displayed. The notation here fixes both boundary components
pointwise, and differs from a convention allowing their permutation.

Gluing an annulus between these boundaries sends both $\delta_i$ to $t$.
Cutting a representative of a mapping class preserving $\vec\gamma$
produces a preimage. Its ambiguity is the opposite pair of boundary
twists $\delta_1\delta_2^{-1}$. Therefore the discrete sewing map has
kernel generated by $(z_4,-2)$, with an immaterial inversion of this
generator. Put
\[
 E=(\widehat B_4\times\widehat{\mathbf Z})/
       \overline{\langle(z_4,-2)\rangle},\qquad t_E=[(1,1)].
\]
The closed subgroup by which we divide is central. There is a natural
continuous sewing map $E\to(\widehat\Gamma(S))_{\vec\gamma}$.
This map is an isomorphism, including its central inertia, as follows.

First, $a\mapsto t_E^a$ is an injection of $\widehat{\mathbf Z}$.
The continuous homomorphism
\[
 (b,a)\longmapsto\operatorname{exp}(b)+6a
      \quad:\quad\widehat B_4\times\widehat{\mathbf Z}
           \longrightarrow\widehat{\mathbf Z}
\]
vanishes on $(z_4,-2)$, since $\operatorname{exp}(z_4)=12$.
It sends $t_E^a$ to $6a$. Multiplication by $6$ on
$\widehat{\mathbf Z}=\prod_p\mathbf Z_p$ is injective.
Second, quotienting $E$ by this inertia gives
\[
 E/t_E^{\widehat{\mathbf Z}}
   =\widehat B_4/\overline{\langle z_4\rangle}
   \cong\widehat{P\Gamma}_{1,2}.
\]
The last identification follows from the same boundary presentation by
capping the two boundaries by two individually labelled punctures.
Completion commutes with this specified quotient: finite quotients on
both sides are exactly the finite quotients of $B_4$ killing $z_4$.

For the intrinsic stabilizer, the opened source is Boggi,
\emph{Automorphisms of profinite mapping class groups}, arXiv
2011.15075v6, Theorem 2.3(i). In the present one-curve case its second
exact sequence reads
\[
 1\longrightarrow t^{\widehat{\mathbf Z}}
 \longrightarrow(\check\Gamma(S))_{\vec\gamma}
 \longrightarrow P\check\Gamma_{1,2}\longrightarrow1.
\]
The two newly created punctures are fixed individually, precisely
because the curve is oriented. Theorem 2.3 states that the left-hand
group is the profinite completion of the twist group. Proposition
2.6(i) states, in its own notation,
\[
 \check\Gamma(S,\partial S)=\widehat\Gamma(S,\partial S)
       \quad\text{for }g(S)\leq2,
\]
and the same equality for the pure groups. The paragraph immediately
preceding Theorem 2.3 also gives the genus-at-most-two comparison for
punctured surfaces. Thus the checks may be replaced by hats in the
displayed exact sequence. This is the exact profinite input needed;
no assertion about all-genus congruence is used.

The sewing map between these two exact sequences is the identity on
$\widehat{\mathbf Z}$ and on $\widehat{P\Gamma}_{1,2}$.
For completeness, if an element of its kernel is given, its image in
the quotient is trivial; hence it is $t_E^a$, and the equality on
inertia forces $a=0$. Conversely, for an element of the target, lift
its quotient to $E$. The difference lies in target inertia and is
also in the image. The map is therefore bijective and, being a
continuous bijection between profinite groups, is an isomorphism.
This proves the first line of E.64.3 with $t_E^2=z_4$.

For its action formula, use the standard positive genus-two chain
$T_1,\ldots,T_5$, put $t=T_5$ and identify the indicated $B_4$ with
$\langle T_1,T_2,T_3\rangle$. The three-chain relation gives
$(T_1T_2T_3)^4=T_5^2$. The element
$z_5=(T_1T_2T_3T_4)^5$ is the hyperelliptic involution: its sphere
quotient is trivial (the complementary disc contains only one marked
point), whereas its homological action is $-I$. In particular it is
central. Consequently $f(z_5,T_5^2)=1$ for every commutator word $f$.
The standard GT formula therefore sends $T_5$ to $T_5^\lambda$ and
restricts on the first three generators to $\Phi_{4,F}$.
On $E$ this is exactly
\[
 [(b,d^a)]\longmapsto[(\Phi_{4,F}(b),d^{\lambda a})].
\]
It respects $(z_4,-2)$ since $\Phi_{4,F}(z_4)=z_4^\lambda$.
Here the descent of the standard braid action to the genus-two
hyperelliptic quotient is the separate presentation argument in
E.64.2. This verifies its claimed \emph{selected} sewing square.
It does not identify an arbitrary orbifold action with a restriction
of this action.

\begin{statement}{B.27}{Known}{E.64.13}
\citep[Propositions 4.11 and 4.13]{LitNakamuraTsunogaiYasuda2010}
\end{statement}
Use exactly the braid generators $(a,b,c)$ and Ferrari basis
\[
 \xi_1=ca^{-1},\qquad
 \xi_2=aba c^{-1}b^{-1}a^{-1}
\]
of E.64. Let $\theta:\widehat B_3^{\rm mor}\to\widehat B_4$
be its specific subgroup map with $\theta(u)=ab$. For
$\sigma\in G_{\mathbf Q}$ assume $\chi(\sigma)=1$.
Let $A_\sigma$ denote its positive pure order-three frame and
$\Omega_{3,\sigma}$ the original triangle action, lifted as in
E.64.28. The source statements are Nakamura--Tsunogai--Yasuda
(2010), Proposition 4.11 (p.444) and Proposition 4.13 (p.445).
Item 4.12 is a subsection, not a proposition.

In the source's notation $a=\tau_1$, $b=\tau_2$, $c=\tau_3$.
Proposition 4.11 gives at multiplier one
\[
 u_\sigma=h_\sigma(a^2,b^2)^{-1}(ab)^{3\rho_K(\sigma)}.
\]
The $\lambda=1$ specialization of the source's (1.1) is
$f(x,y)=h(y,z)^{-1}h(x,y)$. Its conjugacy equation identifies
$h_\sigma(a^2,b^2)^{-1}$ with $A_\sigma$ modulo the three-strand
centre; both lie in the same free pure complement, so uniqueness
of the positive pure frame gives equality there. Equivalently,
one can compare $f^{-1}(ab)=h^{-1}(ab)h$, which is the order-three
frame equation in this convention. Thus
\[
 u_\sigma=A_\sigma(ab)^{3\rho_K(\sigma)}.
\]
Proposition 4.13 gives, on the three triangle generators,
\[
 \Phi_{4,\sigma}\theta
       =\operatorname{Inn}(u_\sigma)\theta\Omega_{3,\sigma}.
\]
These are equations in the full completed braid group. In the
source proof the initially possible power of its centre is killed
by the exponent sum $\widehat B_4\to\widehat{\mathbf Z}$; the
full twist has exponent $12$, and multiplication by $12$ is
injective on $\widehat{\mathbf Z}$. No residual central sign
is silently discarded.

Normalize by the frame and apply $\theta^{-1}$. The homomorphism
$\theta$ sends $u^{3\rho_K}$ to $(ab)^{3\rho_K}$, so
\[
 \begin{split}
 \Psi_{3,\sigma}
 &=\theta^{-1}\operatorname{Inn}(A_\sigma^{-1})
                      \Phi_{4,\sigma}\theta\\
 &=\operatorname{Inn}(u^{3\rho_K(\sigma)})\Omega_{3,\sigma},
 \end{split}
\]
and hence
\[
 \Psi_{3,\sigma}\Omega_{3,\sigma}^{-1}
     =\operatorname{Inn}(u^{3\rho_K(\sigma)}).
\]
On $\widehat K^{\rm ab}$ the conjugation action of $ab$ is
$R_+=UL=\left(\begin{smallmatrix}0&1\\-1&1\end{smallmatrix}\right)$,
with $R_+^3=-I$. This proves
\[
 \mathcal D_3(\sigma)=(-I)^{\rho_K(\sigma)}.
\]
The central element $(ab)^3$ of the first three-strand subgroup
is distinct from $z_4$. It acts by $-I$ on the Ferrari homology,
whereas $z_4$ acts trivially. This distinction explains the sign.

The parity can be compared to the referee's common radical without
any informal division by two. Definition 4.7 (p.443) fixes
\[
 K^{1/n}=3^{1/(2n)}4^{-1/n}\zeta_{12n}^{-5}.
\]
For $n=2$, with $r_0=(3/4)^{1/4}>0$, this is
$K^{1/2}=(r_0/\sqrt2)\zeta_{24}^{-5}$. The other factors and
$K$ itself lie in $\mathbf Q^{\rm ab}$. Therefore
\[
 (-1)^{\rho_K(\sigma)}=\frac{\sigma(r_0)}{r_0}
       \qquad(\sigma\in G_{\mathbf Q^{\rm ab}}).
\]
This is the same bit used in the repaired Galois part of E.65.

\begin{statement}{B.28}{Open}{the all-GT comparison is not supplied by those sources}
\citep{LitLochakNakamuraSchneps2026,LitNakamuraTsunogaiYasuda2010}
\end{statement}
Proposition 4.4 of the same paper constructs the original triangle
Tate action for every $F\in\widehat{\mathrm{GT}}$. Propositions
4.11 and 4.13, and Theorem 4.16, explicitly have the different
domain $G_{\mathbf Q}$. The introduction, p.432, says:
\begin{quote}
``Unknown is whether they hold for all elements of GT''.
\end{quote}
Here ``they'' refers to the power-conjugate braid formulas of
Proposition 4.13. Theorem 4.16 gives the matrix comparison on
Galois and cannot be extended by replacing its quantifier.

Likewise, Lochak--Nakamura--Schneps (2026), Theorem 1.1, gives
the full profinite mapping-class actions with at most one marked
point. Its Section 8 explicitly leaves the required preservation
of cut-surface images and Lego compatibility for further work.
Boggi's stabilizer theorem identifies the \emph{groups}; it does
not identify two independently constructed actions on them.
The exact missing assertion for A.4(e) is an all-GT computation
of the comparisons $C_3(F)B_{3,F}^{-1}$ and
$C_2(F)B_{4,F}^{-1}$, with the stated paths and normalizations.
Neither the selected sewing square nor the cited full-domain
triangle action makes these quotients trivial.

\begin{statement}{B.29}{Proved}{A convention correction, not a change to the bit}
\citep[Proposition 2.4 and Definition 5.1]{LitNakamuraSchneps2000}
\end{statement}
Nakamura--Schneps (2000), Proposition 2.4 and Definition 5.1,
give the abelianization of their harmonic parameter as
$g_\sigma^{\rm ab}=(xy)^{\rho_2(\sigma)}$, where $\rho_2$
uses the positive roots of $2$. With the E.64 dictionary
$g=\beta^{-1}$ and the E.65 definition of $\rho$ from
$\beta^{-1}$, this is $\rho=\rho_2$, not $-\rho_2$.
The minus sign in the referee report's Marin-convention remark
should therefore be removed after identifying its generators
with the source's generators. The positive shear $+8\rho_2$
in the proof of NTY Theorem 4.16 independently checks this
choice. Its reduction modulo two is unaffected by changing
sign, so the correction changes neither $\delta$ nor $\Delta$.

\subsection{Higher-point rigidity and the manuscript's five-point configurations}


\begin{statement}{B.30}{Refuted}{universal rigidity in six or seven points}
\end{statement}
Pure centralizers of torsion need not be trivial for six or seven marked
points. On the sphere marked at $\{\pm1,\pm2,\pm3\}$, use the involution
$z\mapsto-z$. On the sphere marked at
$\{0,\pm1,\pm2,\pm3\}$, use the same involution. In either case the
circle $|z|=5/2$ is invariant and essential: it separates four from two
punctures in the first case and five from two in the second.
Its Dehn twist is pure, has infinite order, and commutes with the
involution, because conjugation sends a Dehn twist to the twist about the
image curve. Residual finiteness injects the discrete mapping class group
into its profinite completion, so this remains a nontrivial pure
centralizer there. A frame, when it exists, is therefore a torsor under
this centralizer and is not unique.

\begin{statement}{B.31}{Proved}{the count twelve}
\end{statement}
A pentagonal structure on a labelled five-element set is a cyclic order
modulo reversal. There are $5!/10=12$, since its stabilizer is the dihedral
group of order ten. There are six subgroups of order five in
$\mathfrak S_5$; each supports two pentagonal structures, represented by
step one and step two around the same five-cycle. Thus twelve counts
unoriented cyclic orders, not order-five cyclic subgroups. This is exactly
the interpretation of cote 148, manuscript folio 29, scan 31, in the
supplied first-pass transcription.
For the geometric configuration $\{\zeta_5^i\}_{i\in\mathbb Z/5}$,
the arithmetic permutation of the five labels is
$i\mapsto\chi(\sigma)i\pmod5$. Its action on the twelve structures
is the ordinary $\mathfrak S_5$ action on
$\mathfrak S_5/D_5$. It has cycle types
\[
 \begin{array}{c|ccc}
 \chi\bmod5&1&2,3&4\\\hline
 \text{cycle type on twelve structures}&1^{12}&2^2 4^2&1^4 2^4.
 \end{array}
\]
In particular the label permutation $i\mapsto-i$ does not fix all
twelve structures.

There is a different, $\mathfrak S_5$-equivariant operation on this
set: replace a generating five-cycle by its $a$-th power, retaining
its cyclic subgroup. Modulo reversal, this operation fixes all twelve
structures for $a=\pm1$ and interchanges the two structures on each
of the six cyclic subgroups for $a=\pm2$, giving six transpositions.
It is not the same as applying the label permutation $i\mapsto ai$.
The exact program enumerates both actions separately and gives their
different cycle types. The former assertion that these actions coincide
has therefore been corrected.

Lochak--Schneps's two-case rule concerns the two special endpoints
corresponding to the two cocycle classes of A.34, and has the same
$\{\pm1\}$ versus $\{\pm2\}$ dichotomy as the generator-power
operation. That does not identify it with the arithmetic label action
on all twelve cyclic orders. The primitive $k$ describes a path to
the chosen special endpoint; an action on that path space, an action
on the two endpoints, and an action on twelve labelled structures
require separate comparison maps. No such identification is inferred
from their counts alone.

\begin{statement}{B.32}{Proved}{nontrivial automorphism groups of five-point sets}
\end{statement}
For an unordered five-point subset of $\mathbb P^1(\mathbb C)$, its
full nontrivial group of holomorphic automorphisms is one of
$C_2,C_4,D_3,D_5$. To see this, finite subgroups of
$\operatorname{PGL}_2(\mathbb C)$ are cyclic, dihedral, $A_4,S_4,A_5$.
The latter three have no invariant five-element set: their exceptional
orbit sizes are respectively $(4,4,6)$, $(6,8,12)$ and $(12,20,30)$, and
larger orbits cannot help. For a cyclic rotation of order $m$, nonfixed
orbits have size $m$ and there are at most two fixed points; hence
$5=km+r$, $0\leq r\leq2$. For $m>2$ this gives $m=3,4,5$.
The cases $m=3$ and $5$ automatically admit inversion after normalizing
the single nonfixed orbit to $\mu_m$, giving $D_3,D_5$; the case $m=4$
contains one of the two poles and has group $C_4$. For dihedral groups,
exceptional orbit sizes are $2,m,m$, excluding every other case.
The remaining one-dimensional locus has generic group $C_2$.
This classifies smooth configurations with nontrivial symmetry; the
additional vertices of the 542-vertex manuscript diagram include specified
real, boundary or decorated configurations and are not further isolated
smooth automorphism types.

\begin{statement}{B.33}{Open}{the 1620 geometric arrows}
\end{statement}
The supplied transcription gives nine orbit sizes summing to $542$ and
$1620$ arrows. If those arrows are taken as independent oriented edge
generators of a connected graph, its free fundamental group has rank
$1620-542+1=1079$. This is merely the graph calculation, not a presentation
of the geometric fundamental groupoid. Such a presentation additionally
requires the actual homotopy class of each arrow and a proof of the
relations among them. The transcription has explicitly uncertain labels
and the scans are not supplied. No claimed fundamental geometric system
of relations is inferred from these counts.

\begin{statement}{B.34}{Proved}{the genus-two fibre needs another branch point}
\end{statement}
The hyperelliptic curve $y^2=x^5-1$ is branched over
$\mu_5\cup\{\infty\}$, six points, and has genus two by
Riemann--Hurwitz. A sphere with the five marked points $\mu_5$ does not by
itself specify this double cover: a hyperelliptic double cover has an even
number of branch points. Choosing a sixth point, or one of the two fixed
poles of the order-five rotation, is additional data. Therefore the
proposed coordinate in a quotient of $\mathbb Z[\zeta_5]$ and the
$\operatorname{Sp}_4(\mathbb Z/12)$ comparison require a proved
functorial extension to that data; they cannot be read from the bare
five-point notation alone.

\subsection{The quartic comparison proved directly}


\begin{statement}{B.35}{Proved}{the multiplier-one quartic frame without E.39}
\citep[Theorem B$^{\prime}$]{HerfortRibes1985}
\end{statement}
The quartic frame needed in LM7\_PUSH E.64.7 exists uniquely on
$\widehat{GT}\cap\mathcal H$. This assertion does not require the
unverified general periodic-braid statement E.39. We now prove it,
including its exact linear comparison, in the genuine normalization of
E.65.8. In this paragraph only, write
\[
 B_4=\langle a,b,c\mid aba=bab,\ bcb=cbc,\ ac=ca\rangle,
 \qquad d=abc,
\]
so that $d$ is the order-four element called $b$ after passage to $G_5$
in A.33--A.38. Let
\[
 q:\widehat B_4\longrightarrow\widehat B_3,
 \quad q(a)=q(c)=a,\quad q(b)=b,\qquad
 s=aba,\quad z_3=(ab)^3=s^2,\quad z_4=d^4.
\]
Let $j$ be the section generated by the first three strands. Put
\[
 \xi_1=ca^{-1},\qquad
 \xi_2=aba c^{-1}b^{-1}a^{-1},\qquad K=\ker q.
\]
For every $\Phi\in\widehat{GT}\cap\mathcal H$, there is exactly one
$H_\Phi\in\widehat B_4$ such that
\[
 \Phi_4(d)=H_\Phi dH_\Phi^{-1},\qquad q(H_\Phi)=\beta_\Phi.
 \tag{FP.26.1}
\]
Here $\beta_\Phi$ is the unique harmonic pure frame lifted to the free
complement $\langle a^2,b^2\rangle$ in $\widehat B_3$.

\emph{The fold and its free kernel.}
The following exact Nielsen action is obtained from the three Artin
relations:
\[
 \begin{array}{c|cc}
 &\xi_1&\xi_2\\ \hline
 a&\xi_1&\xi_2\xi_1\\
 b&\xi_1\xi_2^{-1}&\xi_2.
 \end{array}
 \tag{FP.26.2}
\]
These two automorphisms of $F(\xi_1,\xi_2)$ satisfy the braid relation.
Conversely, in the resulting semidirect product $F_2\rtimes B_3$, set
$c=\xi_1a$. The elements $a,b,c$ satisfy the three Artin relations:
$a$ fixes $\xi_1$, and
\[
 bcb=\xi_1\xi_2^{-1}aba=cbc,
\]
because $(ab)\xi_1(ab)^{-1}=\xi_2^{-1}$.
Moreover the displayed braid word defining $\xi_2$ gives exactly that
free generator. Thus the two presentations give mutually inverse
homomorphisms, proving $B_4=F(\xi_1,\xi_2)\rtimes B_3$ and
$K=F(\xi_1,\xi_2)$ without importing a Ferrari-kernel assertion.
Every finite characteristic quotient of this finitely generated free
kernel has a finite automorphism group; its action therefore factors
through a finite quotient of $B_3$. These characteristic quotients are
cofinal, so completion gives the exact split extension
\[
 \widehat B_4=\widehat K\rtimes\widehat B_3.
 \tag{FP.26.3}
\]

\emph{Equivariance in the fixed normalization.}
Write $f_0=f(a^2,b^2)$ and $e=f(z_3,c^2)$. The multiplier-one
Drinfeld formulas on these literal generators are
\[
 \Phi_4(a)=a,\qquad
 \Phi_4(b)=f_0^{-1}bf_0,\qquad
 \Phi_4(c)=e^{-1}ce.
 \tag{FP.26.4}
\]
For the last formula the usual first argument is $ba^2b$.
The identity $z_3=a^2(ba^2b)$ permits its replacement by $z_3$:
$a^2$ commutes with both arguments, and $f$ has zero exponent sum in
each variable. Since $q(z_3)=z_3$ commutes with $a^2$,
$q(e)=f(z_3,a^2)=1$. Consequently
\[
 q\Phi_4=\Phi_3q,\qquad \Phi_4j=j\Phi_3.
 \tag{FP.26.5}
\]
This proves preservation of $\widehat K$ as well. Collas (2012),
Section 1, records the standard braid formulas in his opposite
conjugator convention; our convention here is fixed by (FP.26.4).
His prime-order theorem alone does not supply the quartic claim.
The order-four input is instead the argument of A.37--A.38.

\emph{Lifting the five-point frame.}
The quotient $B_4/\langle z_4\rangle$ is the label-five stabilizer in
$\Gamma_{0,[5]}$. Its completion is the corresponding open subgroup
$L$ of $G_5$, and the central kernel is exactly
$\overline{\langle z_4\rangle}\cong\widehat{\mathbb Z}$.
The last isomorphism follows also from exponent sum, which takes
$z_4$ to $12$: the induced topology on this infinite cyclic group is
its full profinite topology. By A.38 there is a pure element
$\bar H\in P_5\subset L$ conjugating $\bar d$ to
$\Phi(\bar d)$. Lift it to $H_0\in\widehat B_4$. For some
$t\in\widehat{\mathbb Z}$ one has
\[
 \Phi_4(d)=H_0dH_0^{-1}z_4^t.
\]
The exponent-sum homomorphism $\widehat B_4\to\widehat{\mathbb Z}$
is fixed by a multiplier-one action. The two sides have exponent sums
$3$ and $3+12t$, so $12t=0$; multiplication by $12$ is injective on
$\widehat{\mathbb Z}$, hence $t=0$. This proves exact conjugacy
upstairs, including the central power which a projective argument alone
would leave undetermined.

The harmonic frame also lifts exactly to $\widehat B_3$:
its possible central discrepancy is killed by the same exponent-sum
argument, now with $z_3$ of exponent sum six. Applying (FP.26.5) gives
\[
 q(H_0)sq(H_0)^{-1}=\beta_\Phi s\beta_\Phi^{-1}.
\]
The centralizer of $s$ in $\widehat B_3$ is
$\overline{\langle s\rangle}\cong\widehat{\mathbb Z}$.
Indeed its image in $C_2\amalg C_3$ lies in the order-two factor by
the free-profinite-product centralizer theorem; its remaining factor
is a central power of $z_3=s^2$. Therefore
$q(H_0)=\beta_\Phi s^a$ for a unique $a\in\widehat{\mathbb Z}$.
Replacing $H_0$ by $H_0d^{-a}$ proves existence in (FP.26.1).

\emph{Uniqueness and the quartic group.}
Equations (FP.26.2) give
\[
 d=\xi_2^{-1}s,\qquad
 d\xi_1d^{-1}=\xi_2^{-1},\qquad
 d\xi_2d^{-1}=\xi_1.
 \tag{FP.26.6}
\]
The preimage of $\overline{\langle s\rangle}$ is consequently
$\widehat K\rtimes\overline{\langle d\rangle}$, with the displayed
order-four action. If $y=b^{-1}d^2b$, direct braid multiplication gives
$yd^{-2}=\xi_2^{-1}$ and $y^2=d^4$. With
$k_0=\xi_2^{-1}$, $k_1=\xi_1^{-1}$, the action is
$k_0\mapsto k_1$, $k_1\mapsto k_0^{-1}$.
Thus the mutually inverse presentation maps
\[
 \langle D,Y\mid Y^2=D^4\rangle
 \ \rightleftarrows\ F(k_0,k_1)\rtimes\mathbb Z,
 \quad D\longmapsto d,\quad Y\longmapsto k_0d^2
\]
identify this preimage with the profinite completion of the quartic
group. Its quotient by $\overline{\langle d^4\rangle}$ is
$C_4\amalg C_2$. The image of $\widehat K$ is the kernel of the
map to $C_4$ that sends $D$ to $1$ and $Y$ to $2$.
If $k\in\widehat K$ centralizes $d$, its image in this free profinite
product lies in the factor $\langle\bar d\rangle$ and also in that
kernel, so it is trivial. Then $k\in\overline{\langle d^4\rangle}$,
and $q(k)=1$ forces $k=1$, since $q(d^4)=s^4$ has infinite
profinite order. Hence
\[
 C_{\widehat K}(d)=1.
 \tag{FP.26.7}
\]
Two normalized frames have ratio in this centralizer, proving uniqueness.
Alternatively A.36, followed by the fold, gives the same conclusion.
The normalized frame is continuous: the map
$H\mapsto(HdH^{-1},q(H))$ is a continuous injection from a compact
space onto its image after (FP.26.7).

Uniqueness and (FP.26.5) also give the exact crossed law
\[
 H_{\Phi\Psi}=\Phi_4(H_\Psi)H_\Phi.
\]
Consequently $\operatorname{Inn}(H_\Phi^{-1})\Phi_4$ is an actual
action on the quartic subgroup, fixing $d$ and $z_4$.

\emph{The signed linear comparison.}
On $\widehat K^{\mathrm{ab}}$ in the ordered basis
$(\xi_1,\xi_2)$, (FP.26.2) and (FP.26.6) give
\[
 Q(a)=U=\begin{pmatrix}1&1\\0&1\end{pmatrix},\qquad
 Q(b)=L_0=\begin{pmatrix}1&0\\-1&1\end{pmatrix},\qquad
 Q(d)=S=UL_0U.
\]
Write $M_\Phi$ for the action of $\Phi_4$ on this homology.
Both arguments of $e$ commute with $a$, so (FP.26.4) gives
$\Phi_4(\xi_1)=e^{-1}\xi_1e$ with $e\in\widehat K$.
Thus $M_\Phi e_1=e_1$. Covariance at $a$ gives
$M_\Phi U=UM_\Phi$, which implies
\[
 M_\Phi=\begin{pmatrix}1&t_\Phi\\0&1\end{pmatrix},\qquad
 M_\Phi L_0M_\Phi^{-1}
       =Q(f_0)^{-1}L_0Q(f_0).
\]
Because $q(H_\Phi)=\beta_\Phi$, (FP.26.3) writes it uniquely as
$H_\Phi=j(\beta_\Phi)k_\Phi$ with $k_\Phi\in\widehat K$.
Inner conjugation by $k_\Phi$ acts trivially on
$\widehat K^{\mathrm{ab}}$. The normalized quartic action therefore
has the literal matrix
\[
 \boxed{\ C_2(\Phi)=Q(\beta_\Phi)^{-1}M_\Phi\ },\qquad
 C_2(\Phi)S=SC_2(\Phi).
 \tag{FP.26.8}
\]
No sewing theorem, unspecified path, or assumed quartic instance of E.39
enters this identity. It proves the quartic input used in E.64 at
multiplier one. It does not by itself prove E.39 for arbitrary numbers
of strands or arbitrary multipliers, nor identify the normalized quartic
action with a separate orbifold action of the same GT pair.
The discrete identities and the integral homology signs are independently
verified by the faithful Artin representation in
\texttt{check\_quartic\_frame\_input.py}; its complete output is retained.

\subsection{The complete checkable E.67 finite calculation}


\begin{statement}{B.36}{Proved}{The complete common correction image and the full commutator image}
Let $F=\widehat F(X,Y)$ and let $q_{12}$ send $X$ to
$U^2$ and $Y$ to $L^2$ in $\mathrm{SL}_2(\mathbb Z/12)$.
Put $N=\ker q_{12}$ and $K_6=\overline{N'N^6}$.
Let $V=\mathbb F_3^2$, and form the two affine groups
\[
 \begin{array}{c|ccc}
 &x&y&z\\ \hline
 M&(0,R)&(e_1,R^3)&(e_1,R^2)\\
 G&(e_1,S)&(-e_2,S^2)&(0,S)
 \end{array}
 \quad
 R=\begin{pmatrix}0&1\\-1&1\end{pmatrix},\quad
 S=\begin{pmatrix}0&1\\-1&0\end{pmatrix}.
\]
Multiplication is $(v,A)(w,B)=(v+Aw,AB)$. Evaluate one free
word in the six ordered pairs
\[
 (x,y),(y,x),(z,x),(x,z),(y,z),(z,y)
\]
first in $M$ and then in $G$, obtaining $\Theta$.
Then $|q_{12}(F)|=96$, $N$ has a Schreier free basis of length
$97$, and
\[
 \begin{split}
 \Theta(K_6)=\Theta(N')&=\mathcal M\times\mathcal G,\\
 \mathcal M&=\{(aw,aw,bw,cw,-cw,-bw):a,b,c\in\mathbb F_3\},\\
 \mathcal G&=\{(u,v,t,-t,-Sv,-Su):u,v,t\in V\},
 \qquad w=(1,-1).
 \end{split}                                                     \tag{ER1}
\]
This is a translation space of dimension nine, with projection
dimensions three and six. The full image $\Theta(F')$ is a
translation space of dimension fifteen. Its Mordell projection
$\mathcal F_M$ has dimension nine, its Gaussian projection is
exactly $\mathcal G$, and
\[
                         \Theta(F')=\mathcal F_M\times\mathcal G.
\]
The literal pair $(x_M,x_G),(y_M,y_G)$ generates $M\times G$ of
order $54\cdot36=1944$.
\end{statement}

\begin{proof}
The independent program
\texttt{agent\_literature/code/recompute\_e67\_structure.py}
represents each affine element by its actual permutation of the
nine points of $V$. It imports no supplied checker. It constructs
the $96$-vertex Cayley graph of $U^2,L^2$ using only positive
letters, with the letter $Y$ explored first. Its prefix-closed
spanning tree has $95$ edges. Of the $192$ positive directed
edges, the other $97$ give the nontrivial Schreier words
$t_g a t_{ga}^{-1}$. Each is freely reduced and re-evaluated to
the identity modulo twelve. The Schreier theorem therefore gives
a free basis of $N$; this argument also identifies the full
profinite kernel, since the subgroup has finite index.
The signed and projective kernels agree here: every element of
$q_4(F)$ has diagonal entries one, whereas $-I$ has diagonal
entries three.

The linear parts of every component of $\Theta(F)$ belong to
a cyclic group, so every component of $\Theta(F')$ is a
translation. Consequently the required derived images are
ordinary vector subspaces over $\mathbb F_3$. The image of
$N'$ is the normal closure of the pairwise commutators of the
$97$ Schreier generators. The program evaluates all
$\binom{97}{2}=4656$ commutators as permutations, extracts their
translation vectors, and obtains rank nine. It then closes this
space under conjugation by all $97$ generator images; the rank
remains nine. This proves that the space is exactly $\Theta(N')$,
not merely a lower bound obtained from a few commutators.
Its basis is compared with the nine basis vectors obtained from
the displayed parametrization in (ER1); their joint rank is nine.
The two projection ranks are three and six, proving the asserted
direct product.

For the sixth-power step, every Mordell component of
$\Theta(N)$ has linear part in $\langle R^2\rangle$, and every
Gaussian component has linear part in $\langle S^2\rangle$.
The respective entire affine subgroups have orders $27$ and $18$.
The program checks exponent dividing six on every element of
these two groups, not just on the $97$ sampled images.
Equivalently, $R^2=I+J$ with $J^2=0$ in characteristic three,
so $I+R^2+R^4=0$, while the second linear part is $I$ or $-I$.
Thus $\Theta(N^6)=1$, and $\Theta(K_6)=\Theta(N')$ follows.

To determine the full commutator image, the program starts from
the single actual word $[X,Y]$ and closes its translation vector
under conjugation by $X,Y,X^{-1},Y^{-1}$. This is exactly the
normal closure defining $F'$, so it gives the entire image.
Its stable dimension is fifteen, with projection dimensions
nine and six. Since their sum equals fifteen and the Gaussian
projection is the already computed $\mathcal G$, the asserted
product decomposition follows.

The JSON certificate saves all $97$ Schreier words, nine actual
words in $N'$ whose images form the correction basis, and fifteen
actual words in $F'$ whose images form the full-image basis.
It re-evaluates every saved basis word, so every vector in the
claimed images has a finite-word preimage. Finally, direct closure
of the two literal pairs gives all $1944$ elements of $M\times G$.
The program and its complete textual output are supplied alongside
the JSON certificate.
\end{proof}

\begin{statement}{B.37}{Proved}{The universal conditional fibre statement is zero or four}
For an arbitrary origin word $f_0\in F'$, its attainable
twelve-evaluation vectors in $f_0(K_6\cap F')$ satisfy the following
exact dichotomy. There are $729$ possible Mordell origin cosets in
$\mathcal F_M/\mathcal M$. Exactly nine contain a framed
determinant-one Mordell target, and each of these contains exactly
one. The other $720$ contain no such target. All four framed
Gaussian targets are reachable from every Gaussian origin.
Consequently an origin in one of the nine Mordell target cosets
has exactly four simultaneous framed targets, whereas every
other origin has none.

Whenever such a target $t$ exists, an actual correction
$c_t\in N'$ can be chosen, and its complete word fibre is
\[
 f_0c_t\bigl(K_6\cap F'\cap\ker\Theta\bigr).                  \tag{ER2}
\]
Each target induces an automorphism of the actual joint group
$M\times G$, hence gives a generating pair there.
\end{statement}

\begin{proof}
The relevant centralizers in $\mathrm{SL}_2(\mathbb F_3)$ have
orders six and four. In the Mordell case the three positive-sign
matrices are $I,R^2,R^4$. For each such matrix the translation
frame of $y$ is unique, because $I-R^3=2I$ is invertible; the
translation frame of $z$ has three possibilities, because
$I-R^2$ has rank one. This gives nine framed Mordell targets.
The cusp condition $(B-I)e_1\in\operatorname{im}(I-R^2)$ forces
$B$ to act as $+1$ on $\mathbb F_3(1,-1)$; all six centralizer
matrices are checked, and the three that pass are exactly those
just listed.

In the Gaussian case the centralizer is $\{I,-I,S,-S\}$.
The map with linear action $B$ on the translation subgroup and
fixing $x$ is conjugation in the full affine permutation group by
$(h,B)$, where
$(I-S)h=(I-B)e_1$. Both $I-S$ and $I-S^2$ are invertible.
Hence both frames are unique and there are exactly four targets.
The first free-word evaluation has translation
$v=(I-B)e_1$, giving
\[
 \begin{array}{c|c|c}
 B&v&\text{parity}\\\hline
 I&(0,0)&0\\-I&(2,0)&0\\S&(1,1)&1\\-S&(1,2)&1.
 \end{array}
\]
Thus the parity is $v_2^2$. These factor actions are implemented
as conjugations of actual nine-point permutations, and the program
verifies that they carry the full finite group bijectively to itself.

If $h_y,h_z$ are the two translation frames, the two-cycle and
three-cycle equations give the six translation values
\[
 (-h_y,h_y,h_z,-h_z,h_y-h_z,h_z-h_y).                            \tag{ER3}
\]
Conversely this tuple satisfies all six ordered two-cycle and
three-cycle evaluations and gives the indicated conjugates of
$y,z$. The program verifies every ordered form for each target.
All nine Mordell target vectors belong to $\mathcal F_M$, and all
four Gaussian target vectors belong to $\mathcal G$.

The quotient $\mathcal F_M/\mathcal M$ has dimension six.
Using the actual echelon normal form for the computed correction
space, the program enumerates its $729$ cosets, verifies that every
coset has $27$ elements, and computes the target distribution
\[
       \#\{\text{cosets with zero targets}\}=720,\qquad
       \#\{\text{cosets with one target}\}=9.
\]
This computation is independent of any list of origin words.
Since the Gaussian correction space equals the entire Gaussian
origin space, all four Gaussian targets are reachable from every
origin. The product decomposition in B.36 makes the two
corrections independent, proving the zero-or-four dichotomy.

Moreover $N'\subseteq K_6\cap F'$ and
$\Theta(N')=\Theta(K_6)$, so restriction to $K_6\cap F'$ loses
no correction vector. A linear combination of the nine saved
word witnesses gives an actual $c_t\in N'$ for every reachable
target. The ordinary kernel-fibre identity for the homomorphism
$\Theta$ gives (ER2). Both factor maps are automorphisms and the
literal pair generates their full direct product, by B.36;
therefore the corresponding joint pair is generating.
\end{proof}

\begin{statement}{B.38}{Open}{The forty-eight original membership tests require their actual words}
The attached material does not include the literal $24$ E.1
origin words and the separate $24$ E.17 origin words needed to
decide their membership in the nine successful Mordell cosets.
This recomputation therefore does not certify that each of those
$48$ specific origins has four targets, nor reproduce their
claimed saved corrections. It proves the complete structural
criterion B.37 and the actual correction image B.36.
For each missing origin, the smallest remaining calculation is
one twelve-coordinate evaluation in $M$, followed by reduction
modulo the three-dimensional space $\mathcal M$ and comparison
with the nine saved target normal forms.
\end{statement}
No failure of an original origin is inferred from the $720$ empty
generic cosets. Conversely, the existence of the nine successful
cosets does not prove membership for a word that has not been supplied.
These are finite necessary comparisons; none supplies a full-GT
lift or a value of $\Delta$ on the pentagon subgroup.

\subsection{Remaining higher-point and character determinations}


\begin{statement}{B.39}{Open}{The corrected six- and seven-point problem}
The complete torsion, centralizer-torsor, special-fibre and
character classification requested in task B.1 has not been
determined. The universal unique-frame premise is false at
both six and seven points, by the explicit invariant essential
twists of B.30.
\end{statement}
A corrected theory must retain a frame as a torsor under its
actual pure centralizer, which in these examples is infinite.
Choosing one element of that torsor does not define a canonical
character. Neither the three rigid five-point classes nor the
arithmetic intersections A.23 give the missing higher-point
classification. The first undecided extension in this return
is already a coordinate independent of the pure-centralizer
choice for those two involution examples. No list of new
six- or seven-point characters is asserted.

\begin{statement}{B.40}{Open}{All relative quadratic characters of the gentle group}
For $K_G$ of P.12, we have not determined
$\operatorname{Hom}_{\rm cont}(K_G,\Ftwo)$, or decided
whether it is generated by $\Delta$.
\end{statement}
B.1 completely classifies the characters that factor through
the proved finite frame pair: they are zero and $\Delta$.
That finite factorization has not been proved for every
relative character. A.81 constructs the actual relative
characters $\Delta\circ D_P^n$ and proves that $\Delta$
is outside their finite span, but does not prove a positive
iterate nonzero. Therefore it does not yet furnish a second
independent quadratic character. The cubic character on
$\HH$ in A.77 is a different object: conjugation by complex
conjugation negates its values, so it is not being counted
as a quadratic character or as an ordinary full-group
relative abelian cubic character.

\begin{statement}{B.41}{Open}{The remaining finite-coefficient and relative-stage homology comparison}
The Alexander and characteristic-zero level-twelve questions
are now resolved in the stated scopes. A.94 excludes a
$\Delta$-line from the canonical Alexander module, and from
homology of stable finite metabelian covers, in characteristic
different from two. A.98--A.99 and A.110 construct a canonical
finite homology action detecting the full character, with a
sign line in its ordered-basis permutation module.
A.111 excludes every $\Delta$ constituent from the raw
level-twelve homology in characteristic zero for the full
gentle group. We have not determined all finite-coefficient
constituents of that raw module or completed the precise
identification with every listed stage of LM8 B.31--B.35.
\end{statement}
The distinction between detecting an action character and
containing it as a raw homology constituent is essential.
A.111 proves the latter nonexistence from an actual
good-reduction Frobenius with no eigenvalue one, using the
genuine arithmetic cover attached to the normalized action.
It does not apply after restricting to $\mathcal H$, and its
determinant need not be invertible in finite coefficients.
The first remaining case is consequently an odd finite
coefficient constituent of the raw level-twelve homology,
with the arithmetic subgroup and full gentle action retained.
B.4 separately
refutes the universal assertion that a module and its
$\Delta$-twist must be nonisomorphic, using
$1\oplus\Delta$. The Galois restrictions are isomorphic
for every module by the proved Galois vanishing of $\Delta$.
